\documentclass{standalone}
\usepackage{circuitikz}
\usetikzlibrary{calc}
\definecolor{circuitnotemark}{HTML}{FE00FE}
\begin{document}
\begin{circuitikz}[american, line width=0.8pt]
\ctikzset{bipoles/length=1.2cm}
\foreach \x in {0,1,2,3} {\foreach \y in {0,-1,-2,-3,-4} {\fill[gray, opacity=0.35] (\x,\y) circle (0.55pt);}}
\node[gray, font=\scriptsize] at (0,0.42) {1};
\node[gray, font=\scriptsize] at (1,0.42) {2};
\node[gray, font=\scriptsize] at (2,0.42) {3};
\node[gray, font=\scriptsize] at (3,0.42) {4};
\node[gray, font=\scriptsize, anchor=east] at (-0.42,0) {a};
\node[gray, font=\scriptsize, anchor=east] at (-0.42,-1) {b};
\node[gray, font=\scriptsize, anchor=east] at (-0.42,-2) {c};
\node[gray, font=\scriptsize, anchor=east] at (-0.42,-3) {d};
\node[gray, font=\scriptsize, anchor=east] at (-0.42,-4) {e};
\coordinate (c2) at (1,-2);
\coordinate (a1) at (0,0);
\coordinate (e1) at (0,-4);
\coordinate (e4) at (3,-4);
\coordinate (a4) at (3,0);
\node[dipchip, num pins=8, font=\scriptsize] (part-U1) at (c2) {$\mathrm{NE555}$}; % line 3
\draw (a1) |- (part-U1.pin 1); % line 5
\draw (e1) |- (part-U1.pin 4); % line 6
\draw (part-U1.pin 5) -| (e4); % line 7
\draw (part-U1.pin 8) -| (a4); % line 8
\path (0,-5.593) rectangle (4.318,-4.407); % line 10
\node[anchor=west, inner sep=0, circuitnotemark, font=\large] at (0,-5) {X}; % line 10
\node[anchor=south west, inner sep=2pt, circuitnotemark, font=\large\bfseries] (circuittitle) at (current bounding box.north west) {X};
\path (circuittitle.south west) rectangle ++(5.8,0.593);
\end{circuitikz}
\end{document}
